\section{Implementierung}
\label{sec:implementierung}

\subsection*{Gesamtaufbau}

Die Webanwendung ist mit Elm~0.19.1 implementiert und erzeugt alle drei
Visualisierungen als SVG. Sie folgt The Elm Architecture mit den Schritten
Model, View und Update \cite{elmArchitecture}. \Cref{fig:architecture} zeigt
den Daten- und Zustandsfluss.

\begin{figure}[htbp]
  \centering
  \begin{tikzpicture}[
    node distance=11mm and 13mm,
    box/.style={draw=projectblue,rounded corners,fill=blue!3,minimum width=35mm,minimum height=9mm,align=center},
    arrow/.style={-{Latex[length=2.5mm]},thick,draw=projectblue}
  ]
    \node[box] (json) {\code{public/data/energy.json}\\744 Stunden per HTTP};
    \node[box,below=of json] (api) {\code{Api.elm}\\Decoder und Fehlerzustand};
    \node[box,below=of api] (main) {\code{Main.elm}\\Model, Update, Zeitfenster};
    \node[box,below left=of main] (network) {\code{FlowNetwork.elm}\\Land wählen};
    \node[box,below=of main] (time) {\code{TimeSeries.elm}\\Stunde wählen};
    \node[box,below right=of main] (matrix) {\code{FlowMatrix.elm}\\Land und Stunde wählen};
    \draw[arrow] (json) -- (api);
    \draw[arrow] (api) -- (main);
    \draw[arrow] (main) -- (network);
    \draw[arrow] (main) -- (time);
    \draw[arrow] (main) -- (matrix);
    \draw[arrow,bend left=18] (network) to (main);
    \draw[arrow,bend right=18] (time) to (main);
    \draw[arrow,bend right=28] (matrix) to (main);
  \end{tikzpicture}
  \caption{Modulstruktur und gerichteter Datenfluss der Elm-Anwendung.}
  \label{fig:architecture}
\end{figure}

Das Modul \code{Domain.elm} enthält die fachlichen Typen. \code{Api.elm}
lädt und decodiert den normalisierten Datensatz per HTTP. Die Views sind
zustandslos: Sie erhalten Daten und Callback-Funktionen, erzeugen SVG und senden
bei Interaktion eine Nachricht an \code{Main.elm}. Dadurch können die drei
Ansichten nicht unbemerkt auseinanderlaufen.

\subsection*{Wichtigste Elm-Datenstrukturen}

Der globale Zustand nach erfolgreichem Laden lautet verkürzt:

\begin{lstlisting}[language={}]
type alias State =
    { dataset : Dataset
    , selectedIndex : Int
    , selectedPartner : Maybe String
    , windowStart : Int
    , windowSize : Int
    }
\end{lstlisting}

\code{selectedIndex} ist ein globaler Index in den 744 Stunden. Für jede View
wird daraus ein lokaler Index des sichtbaren Fensters gebildet.
\code{selectedPartner} ist optional, weil die Erzeugungszeitreihe auch ohne
Länderfilter sinnvoll bleibt. \code{windowStart} und \code{windowSize}
definieren 48-Stunden-, 7-Tage- oder Monatsansicht.

Die wichtigsten Nachrichten sind:

\begin{lstlisting}[language={}]
type Msg
    = GotDataset (Result Http.Error Dataset)
    | SelectPartner String
    | SelectTime Int
    | SelectCell String Int
    | SetWindowSize Int
    | MoveWindow Int
    | Reset
\end{lstlisting}

\code{SelectCell} setzt Partner und Zeit atomar. \code{MoveWindow} begrenzt
den neuen Start mit \code{clamp}; dadurch kann weder vor den ersten noch hinter
den letzten Datensatz gesprungen werden. Ladefehler erzeugen einen sichtbaren
Fehlerzustand statt einer leeren Seite.

\subsection*{SVG-Berechnung}

Der Netzwerkgraph berechnet Positionen auf einem Kreis aus Index und Anzahl der
Partner. Ein quadratischer Bézier-Pfad verbindet jeden Partner mit dem
deutschen Knoten. Das Vorzeichen bestimmt Reihenfolge der Endpunkte,
Pfeilmarker und Farbe; \(1{,}5+\min(20,3{,}2|x|)\) bestimmt die Linienbreite.

Die Zeitreihe summiert für jede Stunde die vier Gruppen. Für jede gestapelte
Fläche wird eine obere und untere Punktfolge zu einem geschlossenen SVG-Pfad
verbunden. Die Y-Skala wird aus dem Maximalwert des sichtbaren Fensters und
einer gerundeten Schrittweite bestimmt. Der Länderfluss besitzt eine getrennte
symmetrische Skala. Über jedem Stundenintervall liegt ein transparenter
Trefferbereich, sodass die Auswahl nicht auf eine dünne Linie beschränkt ist.

Die Pixelmatrix berechnet die Zellbreite aus Fensterlänge und verfügbarer
Breite. Der Betrag wird durch das globale Maximum aller Monatswerte normiert;
danach wird zwischen Hellgrau und Rot beziehungsweise Blau linear interpoliert.
Die Farbnorm hängt somit nicht vom sichtbaren Fenster ab.

\subsection*{Implementierungsaufwand und Entscheidungen}

Relativ direkt war die Umsetzung des Elm-Grundgerüsts, der JSON-Decoder und
einfacher SVG-Elemente auf Basis der Übungen. Aufwändiger waren die
Vorverarbeitung heterogener View-Schemata, die korrekte Einheit von
\code{v\_totalpower}, das Paginieren von mehr als 1000 Rohzeilen, gestapelte
Flächenpfade, zwei quantitative Y-Skalen und die Umrechnung zwischen globalem
und lokalem Zeitindex.

Eine direkte Authentifizierung aus Elm wurde verworfen. Das Passwort müsste in
einer öffentlichen Pages-Anwendung eingegeben oder der Token sichtbar werden.
Stattdessen trennt die Pipeline geschützten Export und öffentliche Analyse. Die
fertige JSON-Datei wird dennoch per HTTP geladen und kann unabhängig vom
Entwicklungsrechner ausgeliefert werden.

\subsection*{Prüfung und Bereitstellung}

Der Produktionsbuild wurde mit folgendem Befehl erfolgreich erzeugt:

\begin{lstlisting}[language={}]
elm make src/Main.elm --optimize --output=public/elm.js
\end{lstlisting}

Ein automatisches Prüfskript validiert Datenumfang, lückenlose Stundenschritte,
Partnerabdeckung und Wertebereiche. Im Browser wurden alle Auswahlpfade sowie
Zeitraumwechsel manuell ausgeführt. Neue JavaScript- oder HTTP-Fehler traten im
optimierten Build nicht auf. Ein reproduzierbares Capture-Skript wiederholt
definierte Klickfolgen über das Browser-Debug-Protokoll. So entstehen die
Berichtsabbildungen direkt aus der laufenden Elm-Anwendung.

Die Datei \code{.gitlab-ci.yml} verwendet ein Node-Abbild, installiert die
fixierte Elm-Version~0.19.1-6, kompiliert mit \code{--optimize} und publiziert
den Ordner \code{public} über GitLab Pages. Quellcode und Anwendung liegen im
\projectlink{öffentlichen Hochschul-GitLab-Projekt}. Die lokale Git-Historie
ist im Anhang dokumentiert.

\subsection*{Einsatz von KI}

Für Entwürfe von Elm-Code, Python-Hilfsskripten, Dokumentation und Bericht wurde
ausschlie{\ss}lich OpenAI Codex verwendet. Themenwahl, Rückfragen, fachliche
Prioritäten und Abnahme lagen beim Verfasser. KI-Ausgaben wurden nicht ungeprüft
übernommen: Der Code wurde kompiliert, Daten wurden mit unabhängigen
Validierungsregeln geprüft, Interaktionen im Browser ausgeführt und jede Seite
des erzeugten PDF visuell kontrolliert. Fehler und unklare Annahmen führten zu
weiteren Prompts und Änderungen. Die relevanten Nutzereingaben sind nach
Programmteil in \cref{app:ki} aufgeführt; KI-Ausgaben werden entsprechend der
Vorgabe nicht wiederholt.
